\documentclass[a4paper]{ltxdoc}
\usepackage[UTF8, heading=true, scheme=plain, linespread=1.2, zihao=-4, fontset=fandol]{ctex}
\usepackage{geometry, parskip, seqsplit, hyperref, fancyhdr, etoolbox, tocloft}
\ctexset{
	space=auto,
	section = {
		format = \centering \heiti \sanhao \bfseries,
		aftername = \hspace{0.5em},
		afterindent = true,
	},
	subsection = {
		format = \heiti \bfseries,
		aftername = \hspace{0.4em},
		afterindent = true,
	},
	subsubsection = {
		format = \songti \bfseries,
		aftername = \hspace{0.4em},
		afterindent = true,
	},
}

\geometry{left=3cm, right=3cm, top=2.54cm, bottom=2.54cm}
\setmainfont{Times New Roman}

% 字号设置
\newcommand{\chuhao}{\fontsize{42pt}{\baselineskip}\selectfont}     % 字号设置
\newcommand{\xiaochuhao}{\fontsize{36pt}{\baselineskip}\selectfont} % 字号设置
\newcommand{\yichu}{\fontsize{32pt}{\baselineskip}\selectfont}      % 字号设置
\newcommand{\yihao}{\fontsize{28pt}{\baselineskip}\selectfont}      % 字号设置
\newcommand{\erhao}{\fontsize{21pt}{\baselineskip}\selectfont}      % 字号设置
\newcommand{\xiaoerhao}{\fontsize{18pt}{\baselineskip}\selectfont}  % 字号设置
\newcommand{\sanhao}{\fontsize{15.75pt}{\baselineskip}\selectfont}  % 字号设置
\newcommand{\sihao}{\fontsize{14pt}{\baselineskip}\selectfont}      % 字号设置
\newcommand{\xiaosihao}{\fontsize{12pt}{\baselineskip}\selectfont}  % 字号设置
\newcommand{\wuhao}{\fontsize{10.5pt}{\baselineskip}\selectfont}    % 字号设置
\newcommand{\xiaowuhao}{\fontsize{9pt}{\baselineskip}\selectfont}   % 字号设置
\newcommand{\liuhao}{\fontsize{7.875pt}{\baselineskip}\selectfont}  % 字号设置
\newcommand{\qihao}{\fontsize{5.25pt}{\baselineskip}\selectfont}    % 字号设置

% 下划线
\newcommand{\underlineFixlen}[2][3.5cm]{\underline{\makebox[#1][c]{#2}}}

\newcommand{\keywords}[1][ＸＸＸ；ＸＸＸ；ＸＸＸ；ＸＸＸ；ＸＸＸ]{#1}
\newcommand{\keywordsEn}[1][×××; ×××; ×××; ×××; ×××]{#1}

% 中文摘要
\renewenvironment{abstract}{
\thispagestyle{empty} % 去掉页码
{
\begin{center}
\Large \songti \bfseries 摘\hspace{1em}要\vspace{1.1cm}
\end{center}
}
\setlength{\parindent}{2em}
\setlength{\parskip}{0em}
\setlength{\baselineskip}{22pt} % （宋体，小四；固定行距22磅，段前、段后均为0行间距。段落首行缩进2字符。）
\songti
}{
\setlength{\parindent}{0em}
\setlength{\parskip}{1em}
{\par \songti \bfseries{关键词：}}
\keywords
\clearpage
}

% 英文摘要
\newenvironment{abstractEn}{
\thispagestyle{empty} % 去掉页码
{
\begin{center}
\Large \bfseries ABSTRACT\vspace{1.5cm}
\end{center}
}
\setlength{\parindent}{1em}
\setlength{\parskip}{0em}
\setlength{\baselineskip}{22pt} % 22磅行距，首行缩进1字符，段前、段后均为0行间距
}{
\setlength{\parindent}{0em}
\setlength{\parskip}{1em}
{\par \bfseries{KEYWORDS:}}
\keywordsEn
\clearpage
}

% 目录
\renewcommand\contentsname{
{
\begin{center}
\songti \Large \bfseries 目\hspace{1em}录 % （宋体，小二号，加粗；居中，单倍行距，段前0.5行、段后1.5行间距）
\end{center}
}
}



\begin{document}
\begin{titlepage}
{\heiti 学\hspace{1.5em}号：\underlineFixlen[3.5cm]{\studentID} \hfill
	{\heiti 密\hspace{1.5em}级：\underlineFixlen[3.5cm]{\privacy}}}

\centering
{\vspace{1.7cm} \includegraphics{images/hfut_name.png}\vspace{0.3cm}}

{\LARGE \bfseries Hefei University of Technology}\vspace{1cm}

{\chuhao \heiti 本科毕业设计（论文）}\vspace{0.7cm}

{\LARGE \bfseries UNDERGRADUATE THESIS}\vspace{0.9cm}

{\includegraphics[width=3.76cm, height=3.76cm]{images/hfut_logo.jpg}\vspace{1.3cm}}

{
\linespread{1.6}
\songti \sanhao
	{\bfseries 类\hspace{2em}型：}\underlineFixlen[8.8cm]{\type}\\
	{\bfseries 题\hspace{2em}目：}\underlineFixlen[8.8cm]{\titleCn}\\
	{\bfseries 专业名称：}\underlineFixlen[8.8cm]{\major}\\
	{\bfseries 入校年份：}\underlineFixlen[8.8cm]{\enrolmentYear}\\
	{\bfseries 学生姓名：}\underlineFixlen[8.8cm]{\studentNameCn}\\
	{\bfseries 指导教师：}\underlineFixlen[8.8cm]{\supervisor}\\
	{\bfseries 系名称\hspace{1em}：}\underlineFixlen[8.8cm]{\department}\\
	{\bfseries 完成时间：}\underlineFixlen[8.8cm]{\finishedYear 年\finishedMonth 月}\\
}


\end{titlepage}

\begin{titlepage}
\centering

{
\parskip=0.5em
\linespread{1.25}
\LARGE \heiti
合\hspace{1.5em}肥\hspace{1.5em}工\hspace{1.5em}业\hspace{1.5em}大\hspace{1.5em}学\vspace{2.95cm}

\bfseries{本科毕业设计（论文）}\vspace{2cm}

\songti \bfseries{\titleCn} \vspace{6cm}
}

{
\parskip=0.5em \linespread{1.5}
\songti \sanhao
学生姓名：\underlineFixlen[8.8cm]{\studentNameCn}

学生学号：\underlineFixlen[8.8cm]{\studentID}

指导教师：\underlineFixlen[8.8cm]{\supervisor}

专业名称：\underlineFixlen[8.8cm]{\major}

系名称\hspace{1em}：\underlineFixlen[8.8cm]{\department}

\vspace{3.2cm}
\large
\finishedYear 年\finishedMonth 月
}


\end{titlepage}

\begin{titlepage}
\centering
{
\parskip=0pt \linespread{1.25}
\sanhao \bfseries{A Dissertation Submitted for the Degree of Bachelor}\vspace{4.7cm}

\Large \bfseries{\titleEn} \vspace{1.8cm}}

{\sanhao By

\studentNameEn
\vfill
Hefei University of Technology

Hefei, Anhui, P.R.China

\finishedMonth\enspace Month, \finishedYear\enspace Year
\vspace{3cm}
}

\end{titlepage}

\begin{titlepage}
\setlength{\parindent}{2em}
\setlength{\parskip}{0.5em}

{
\begin{center}
\heiti \Large
	\bfseries{毕业设计（论文）独创性声明}\vspace{1.2cm}
\end{center}
}

{
本人郑重声明：所呈交的毕业设计（论文）是本人在指导教师指导下进行独立研究工作所取得的成果。据我所知，除了文中特别加以标注和致谢的内容外，设计（论文）中不包含其他人已经发表或撰写过的研究成果，也不包含为获得\underlineFixlen[3cm]{合肥工业大学}或其他教育机构的学位或证书而使用过的材料。对本文成果做出贡献的个人和集体，本人已在设计（论文）中作了明确的说明，并表示谢意。

毕业设计（论文）中表达的观点纯属作者本人观点，与合肥工业大学无关。\vspace{1cm}

毕业设计（论文）作者签名：\hfill 签名日期：\hspace{2em}年\hspace{2em}月\hspace{2em}日
}

{
\vspace{3cm}
\begin{center}
\heiti \Large
	\bfseries{毕业设计（论文）版权使用授权书}\vspace{1.2cm}
\end{center}
}

{
本学位论文作者完全了解\underlineFixlen[3cm]{合肥工业大学}有关保留、使用毕业设计（论文）的规定，即：除保密期内的涉密设计（论文）外，学校有权保存并向国家有关部门或机构送交设计（论文）的复印件和电子光盘，允许设计（论文）被查阅或借阅。本人授权\underlineFixlen[3cm]{合肥工业大学}可以将本毕业设计（论文）的全部或部分内容编入有关数据库，允许采用影印、缩印或扫描等复制手段保存、汇编毕业设计（论文）。

（保密的毕业设计（论文）在解密后适用本授权书） \vspace{1.5cm}

学位论文作者签名：\hfill 指导教师签名：\hspace{7em}

签名日期：\hspace{2em}年\hspace{2em}月\hspace{2em}日\hfill 签名日期：\hspace{2em}年\hspace{2em}月\hspace{2em}日
}


\end{titlepage}


\begin{abstract}
ＸＸＸＸＸＸＸＸＸＸＸＸＸＸＸＸＸＸＸＸＸＸＸＸＸＸＸＸＸＸＸＸＸＸＸＸＸＸＸＸＸＸＸＸＸＸＸＸＸＸＸＸＸＸＸＸＸＸＸＸＸＸＸＸＸＸＸＸＸＸＸＸＸＸＸＸＸＸＸＸＸＸＸＸＸＸＸＸＸＸＸＸＸＸＸＸＸＸＸＸＸＸＸＸＸＸＸＸＸＸＸＸＸＸＸＸＸＸＸＸＸＸＸＸＸＸＸＸＸＸＸＸＸＸＸＸＸＸＸＸＸＸＸＸＸＸＸＸＸＸＸＸＸＸＸＸＸＸＸＸＸＸＸＸＸＸＸＸＸＸＸＸＸＸＸＸＸＸＸＸＸＸＸＸＸＸＸＸＸＸＸＸＸＸＸＸＸＸＸＸＸＸＸＸＸＸＸＸＸＸＸＸＸＸＸＸＸＸＸＸＸＸＸＸＸＸＸＸＸＸＸＸＸＸＸＸＸＸＸＸＸＸＸＸＸＸＸＸＸＸＸＸＸＸＸＸＸＸＸＸＸＸＸＸＸＸＸＸＸＸＸ。

ＸＸＸＸＸＸＸＸＸＸＸＸＸＸＸＸＸＸＸＸＸＸＸＸＸＸＸＸＸＸＸＸＸＸＸＸＸＸＸＸＸＸＸＸＸＸＸＸＸＸＸＸＸＸＸＸＸＸＸＸＸＸＸＸＸＸＸＸＸＸＸＸＸＸＸＸＸＸＸＸＸＸＸＸＸＸＸＸＸＸＸＸＸＸＸＸＸＸＸＸＸＸＸＸＸＸＸＸＸＸＸＸＸＸＸＸＸＸＸＸＸＸＸＸＸＸＸＸＸＸＸＸＸＸＸＸＸＸＸＸＸＸＸＸＸＸＸＸＸＸＸＸＸＸＸＸＸＸＸＸＸＸＸＸＸＸＸＸＸＸＸＸＸＸＸＸＸＸＸＸＸＸＸＸＸＸＸＸＸＸＸＸＸＸＸＸＸＸＸＸＸＸＸＸＸＸＸＸＸＸＸＸＸＸＸＸＸＸＸＸＸＸＸＸＸＸＸＸＸＸＸＸＸＸＸＸＸＸＸＸＸＸＸＸＸＸＸＸＸＸＸＸＸＸＸＸＸＸＸＸＸＸＸＸＸＸＸＸＸＸＸＸＸＸＸＸＸＸＸＸＸＸＸＸＸＸＸＸＸＸＸＸＸＸＸＸＸＸＸＸＸＸＸＸＸＸＸＸＸＸＸＸＸＸＸＸＸＸＸＸＸＸＸＸＸＸＸＸＸＸＸＸＸＸＸＸＸＸＸＸＸＸＸＸＸＸＸＸＸＸＸＸＸＸＸＸＸＸＸＸＸＸＸＸＸＸＸＸＸＸＸＸＸＸＸＸＸＸＸＸＸＸＸＸＸＸＸＸＸＸＸＸＸＸＸＸＸＸＸＸＸＸＸＸＸＸＸＸＸＸＸＸＸＸＸＸＸＸＸＸＸＸＸＸＸＸＸＸＸＸＸＸＸＸＸＸＸＸＸＸＸＸＸＸＸＸＸＸＸＸＸＸＸＸＸＸＸＸＸＸＸＸＸＸＸＸＸＸＸＸＸＸＸＸＸＸＸＸＸＸＸＸＸＸＸＸＸＸＸＸＸＸＸＸＸＸＸＸＸＸＸＸＸＸＸＸＸＸＸＸＸＸＸ。

\end{abstract}

\begin{abstractEn}
\seqsplit{×××××××××××××××××××××××××××××××××××××××××××××××××××××××××××××××××××××××××××××××××××××××××××××××××××××××××××××××××××××××××××××××××××××××××××××××××××××××××××××××××××××××××××××××××××××××××××××××××××××××××××××××××××××××××××××××××××××××××××××××××××××××××××××××××××××××××××××××.}

\seqsplit{×××××××××××××××××××××××××××××××××××××××××××××××××××××××××××××××××××××××××××××××××××××××××××××××××××××××××××××××××××××××××××××××××××××××××××××××××××××××××××××××××××××××××××××××××××××××××××××××××.}
\end{abstractEn}

% 目录
{
\setcounter{page}{0}
\tableofcontents
\thispagestyle{empty} % 去掉页码
}

% 设置页眉
\newgeometry{left=2.8cm, right=2.8cm, top=3cm, bottom=3cm}
\fancyhead{}
\chead{\small 合肥工业大学毕业设计（论文）\vspace{0.3cm}}
\pagestyle{fancy}

% 正文
{
% 重写section指令以便设置section段前1行间距，分页
\pretocmd{\section}{\clearpage \vspace*{-2.0em}}{}{}

\setlength{\parindent}{2em}
\setlength{\parskip}{0.5em}
\setlength{\baselineskip}{22pt}

\section{绪论}
\subsection{ＸＸＸＸＸＸ}
\subsubsection{ＸＸＸＸＸＸ}
ＸＸＸＸＸＸＸＸＸＸＸＸＸＸＸＸＸＸＸＸＸＸＸＸＸＸＸＸＸＸＸＸＸＸＸＸＸＸＸＸＸＸＸＸＸＸＸＸＸＸＸＸＸＸＸＸＸＸＸＸＸＸＸＸＸＸＸＸＸＸＸＸＸＸＸＸＸＸＸＸＸＸＸＸＸＸＸＸＸＸＸＸＸＸＸＸＸＸＸＸ。

ＸＸＸＸＸＸＸＸＸＸＸＸＸＸＸＸＸＸＸＸＸＸＸＸＸＸＸＸＸＸＸＸＸＸＸＸＸＸＸＸＸＸＸＸＸＸＸＸＸＸＸＸＸＸＸＸＸＸＸＸＸＸＸＸＸＸＸＸＸＸＸＸＸＸＸＸＸＸＸＸＸＸＸＸＸＸＸＸＸＸＸＸＸＸＸＸＸＸＸＸ。

\section{ＸＸＸＸＸＸ}
here is footnote.\footnote{footnotes working fine}

}

% 参考文献
{
\clearpage % 分页
\phantomsection % 使得hyperref目录能够跳转到正确的位置
\addcontentsline{toc}{section}{参考文献} % 添加到目录中
\nocite{*} % 添加所有文献
\bibliographystyle{unsrt} % 样式
\bibliography{ref.bib}    % 文献数据库
}

% 致谢
\begin{thank}
本论文是在指导老师ＸＸＸＸＸＸＸＸＸＸＸＸＸＸＸＸＸＸＸＸＸＸＸＸＸＸＸＸＸＸＸＸＸＸＸＸＸＸＸＸＸＸＸＸＸＸＸＸＸＸＸＸＸＸＸＸＸＸＸＸＸＸＸＸＸＸＸＸＸＸＸＸＸＸＸＸＸＸＸＸＸＸＸＸＸＸＸＸＸＸＸ。

ＸＸＸＸＸＸＸＸＸＸＸＸＸＸＸＸＸＸＸＸＸＸＸＸＸＸＸＸＸＸＸＸＸＸＸＸＸＸＸＸＸＸＸＸＸＸＸＸＸＸＸＸＸＸＸＸＸＸＸＸＸＸＸＸＸＸＸＸＸＸ。
\end{thank}

% 附录，必要时
\begin{appendix}
【说明：以下内容可放在附录之内：(1) 正文内过于冗长的公式推导；(2) 方便他人阅读所需的辅助性数学工具或表格；(3) 重复性数据和图表；(4) 论文使用的主要符号的意义和单位；(5) 程序说明和程序全文。可按"附录1  XXX"、"附录2  XXX"、……，分章书写。如无需附录，请删除此页。】
\end{appendix}

\end{document}
